\chapter{Integral Perverse Lattices: the \texorpdfstring{\(p/p^+\)}{p/p+} Formalism}
\label{chap:integral-perverse-lattices}
\label{chap:integral-perverse}
\providecommand{\op}{\mathrm{op}}
\providecommand{\im}{\operatorname{im}}

The pro-etale site removes one artificial difficulty from integral
\(\ell\)-adic cohomology: the coefficient ring \(\Zl\), and more generally the
valuation ring of a finite extension of \(\Ql\), is an honest sheaf.  It does
\emph{not} make integral Verdier duality \(t\)-exact for the usual perverse
\(t\)-structure.  Torsion produces a one-degree defect already on a point.
The correct integral formalism therefore contains two adjacent perverse
\(t\)-structures, denoted \(p\) and \(p^+\).  They become identical after
inverting \(\ell\), and Verdier duality exchanges them before inverting
\(\ell\).

This distinction is not optional bookkeeping.  Ignoring it leads to false
claims that integral intermediate extension is self-dual, that reduction
modulo \(\ell\) is perverse \(t\)-exact, or that a rational perverse sheaf has
a canonical integral model.  This chapter gives the exact replacement for
each of those claims.

\section{Coefficient category and notation}

Let \(\mathcal O\) be the valuation ring of a finite extension \(E/\Ql\), let
\(\varpi\) be a uniformizer, and let \(k_{\mathcal O}=\mathcal O/\varpi\).
The special case \((\mathcal O,E)=(\Zl,\Ql)\) is the one needed for the Weil
conjectures.  For a finite-type scheme \(X\) over a field of characteristic
different from \(\ell\), write
\[
  D_{\mathcal O}(X)=D^b_{\cons}(X_{\proet},\mathcal O)
\]
for the bounded constructible, derived-complete coefficient category of
finite Tor-amplitude constructed in the preceding pro-etale chapters.  The
finite Tor-amplitude condition is retained whenever Verdier biduality is
used.  Put
\[
  K_E:=K\otimes_{\mathcal O}^{\mathrm L}E,\qquad
  K_k:=K\otimes_{\mathcal O}^{\mathrm L}k_{\mathcal O}.
\]
Let
\[
  \D_{X,\mathcal O}K
    :=\R\mathcal Hom_{\mathcal O}(K,\omega_{X,\mathcal O})
\]
be Verdier duality.  On the constructible finite-Tor subcategory it is an
anti-equivalence and the biduality morphism is an isomorphism.

The middle perverse \(t\)-structure \(p\) is defined by the same geometric
stalk and costalk inequalities as over \(E\); see
\cref{chap:perverse-t-structure}.  Its heart is denoted
\(\Perv_p(X,\mathcal O)\), and its cohomology functors by
\({}^p\!H^i\).  An object \(T\) of this heart is \(\varpi\)-power torsion if
\(\varpi^nT=0\) for some \(n\).  It is torsion-free if multiplication by
\(\varpi\) is a monomorphism in the perverse heart.

\begin{warning}[Derived completeness is not the perverse issue]
Derived \(\varpi\)-completeness makes inverse limits and the coefficient
sheaf \(\mathcal O\) behave correctly.  The \(p/p^+\) distinction comes from
the homological dimension-one ring \(\mathcal O\), not from a defect of the
pro-etale site.  Passing from the etale site to the pro-etale site therefore
does not erase this distinction.
\end{warning}

\section{The obstruction is visible on a point}

Consider \(X=\Spec\bar k\).  The perverse \(t\)-structure is the ordinary
cohomological \(t\)-structure.  The free module \(\mathcal O\) is fixed by
\(\R\Hom_{\mathcal O}(-,\mathcal O)\), whereas the resolution
\[
  0\longrightarrow\mathcal O\xrightarrow{\ \varpi\ }\mathcal O
    \longrightarrow k_{\mathcal O}\longrightarrow0
\]
gives
\[
  \R\Hom_{\mathcal O}(k_{\mathcal O},\mathcal O)\simeq
  k_{\mathcal O}[-1].
\tag{\thechapter.1}\label{eq:point-torsion-dual-shift}
\]
Thus duality sends a torsion object in degree \(0\) to a torsion object whose
only cohomology lies in degree \(1\).  No self-dual integral \(t\)-structure
can have both free and torsion modules in its heart in degree \(0\).

The remedy is the following tilted \(t\)-structure on a point:
\begin{align*}
 {}^{+}D^{\leq0}
   &=\{K:H^i(K)=0\ (i>1),\ H^1(K)\text{ is torsion}\},\\
 {}^{+}D^{\geq0}
   &=\{K:H^i(K)=0\ (i<0),\ H^0(K)\text{ is torsion-free}\}.
\end{align*}
Its heart consists of two-term objects with torsion-free \(H^0\) and
torsion \(H^1\).  Formula \eqref{eq:point-torsion-dual-shift} says precisely
that integral duality exchanges the ordinary heart with this tilted heart.

\section{The torsion pair and the \texorpdfstring{\(p^+\)}{p+}-structure}

The full subcategories
\[
\begin{aligned}
  \mathsf T_X
    &=\{P\in\Perv_p(X,\mathcal O):P\text{ is }\varpi\text{-power torsion}\},\\
  \mathsf F_X
    &=\{P\in\Perv_p(X,\mathcal O):\varpi:P\to P\text{ is monic}\}.
\end{aligned}
\]
form a torsion pair in \(\Perv_p(X,\mathcal O)\).  Indeed, the ascending
chain of kernels of \(\varpi^n:P\to P\) stabilizes by constructibility and
Noetherianity.  Its stable value \(P_{\mathrm{tor}}\) is the maximal torsion
subobject, and
\[
  0\longrightarrow P_{\mathrm{tor}}\longrightarrow P
    \longrightarrow P_{\mathrm{tf}}\longrightarrow0
\tag{\thechapter.2}\label{eq:perverse-torsion-decomposition}
\]
has \(P_{\mathrm{tf}}\in\mathsf F_X\).  Moreover
\(\Hom(\mathsf T_X,\mathsf F_X)=0\).

\begin{theorem}[The adjacent integral perverse structure]
\label{thm:p-pplus-tilt}
There is a bounded \(t\)-structure \(p^+\) on \(D_{\mathcal O}(X)\), the
right tilt of \(p\) at \((\mathsf T_X,\mathsf F_X)\), characterized by
\begin{equation}
\begin{aligned}
 {}^{p^+}D^{\leq0}
   &=\{K\in{}^pD^{\leq1}:{}^pH^1(K)\in\mathsf T_X\},\\
 {}^{p^+}D^{\geq0}
   &=\{K\in{}^pD^{\geq0}:{}^pH^0(K)\in\mathsf F_X\}.
\end{aligned}
\label{eq:pplus-tilt-formulas}
\end{equation}
Consequently its heart consists exactly of the objects \(K\) for which
\({}^pH^i(K)=0\) outside \(i=0,1\), with
\({}^pH^0(K)\) torsion-free and \({}^pH^1(K)\) torsion.
\end{theorem}

\begin{proof}
For any torsion pair \((\mathsf T,\mathsf F)\) in the heart of a bounded
\(t\)-structure, the two classes in \eqref{eq:pplus-tilt-formulas} satisfy
orthogonality and provide a truncation triangle.  Explicitly, start with the
\(p\)-truncation triangle of \(K\), split \({}^pH^1(K)\) by
\eqref{eq:perverse-torsion-decomposition}, and take the appropriate cone.
The resulting negative piece has only a torsion \({}^pH^1\); the positive
piece has torsion-free \({}^pH^0\).  This constructs the \(p^+\)-truncation
functors and proves boundedness.

On every smooth stratum the construction is the pointwise tilt above,
shifted by the dimension of the stratum.  Recollement therefore glues the
local tilted structures, and refinement of the stratification does not
change them.  This is the Dedekind-coefficient construction of
\cite[\S3.3.4]{BBD}, now realized inside the pro-etale coefficient category.
\end{proof}

\begin{remark}
The inclusions
\[
  {}^pD^{\leq0}\subset{}^{p^+}D^{\leq0}\subset{}^pD^{\leq1},\qquad
  {}^{p^+}D^{\geq0}\subset{}^pD^{\geq0}\subset{}^{p^+}D^{\geq-1}
\]
make the phrase adjacent \(t\)-structures literal.  Their difference is
entirely \(\varpi\)-primary torsion.
\end{remark}

\section{Verdier duality exchanges \texorpdfstring{\(p\) and \(p^+\)}{p and p+}}

\begin{theorem}[Integral perverse duality]
\label{thm:integral-verdier-exchange}
On constructible objects of finite Tor-amplitude, Verdier duality satisfies
\begin{equation}
\begin{aligned}
  \D_{X,\mathcal O}({}^pD^{\leq0})
    &={} ^{p^+}D^{\geq0},&
  \D_{X,\mathcal O}({}^pD^{\geq0})
    &={} ^{p^+}D^{\leq0},\\
  \D_{X,\mathcal O}({}^{p^+}D^{\leq0})
    &={} ^pD^{\geq0},&
  \D_{X,\mathcal O}({}^{p^+}D^{\geq0})
    &={} ^pD^{\leq0}.
\end{aligned}
\label{eq:duality-exchanges-p-pplus}
\end{equation}
In particular,
\[
  \D_{X,\mathcal O}:\Perv_p(X,\mathcal O)^{\op}\xrightarrow{\sim}
  \Perv_{p^+}(X,\mathcal O).
\tag{\thechapter.5}\label{eq:duality-hearts-p-pplus}
\]
\end{theorem}

\begin{proof}
Choose a finite smooth stratification adapted to \(K\).  Under the
six-operation hypotheses fixed in Chapter~4, constructibility and finite
Tor-amplitude make stalks and costalks dualizable.  Smooth absolute
cohomological purity identifies the dualizing complex on a smooth stratum of
dimension \(d\), with the shift and twist dictated by the standing
normalization.  Thus the support and cosupport inequalities are exchanged by
Verdier duality.  On the coefficient complex the only integral defect is the
point calculation \eqref{eq:point-torsion-dual-shift}: finite free modules
remain in degree \(0\), whereas torsion modules move by one degree.
Verdier duality also exchanges \(j_!\) with \(j_*\).  Induction on the number
of strata and the two recollement triangles therefore glue the assertion
over \(X\); biduality gives the second line from the first.  This is the
middle-perversity case of \cite[\S3.3.3--3.3.4]{BBD}.
\end{proof}

\begin{powerbox}
Over \(E\), the middle perverse structure is self-dual.  Integrally, the
correct statement is not that \(\D\) preserves \(p\), but that \(\D\)
exchanges \(p\) and \(p^+\).  It records exactly where torsion can occur in a
dual complex.
\end{powerbox}

\section{Rationalization collapses the two structures}

Because \(E=\mathcal O[1/\varpi]\) is flat, rationalization is a triangulated
functor and kills precisely the \(\varpi\)-power-torsion constructible
objects.  It is compatible with restriction to geometric points and hence
with all support conditions defining perversity.

\begin{theorem}[Rationalization is \(t\)-exact twice]
\label{thm:rationalization-perverse}
The functor
\[
  (-)_E:D_{\mathcal O}(X)\longrightarrow D^b_c(X_{\proet},E)
\]
is \(t\)-exact from the \(p\)-structure to the middle perverse structure
over \(E\), and also from the \(p^+\)-structure to that same structure.  For
every \(K\) and every \(i\) there are functorial isomorphisms
\[
  ({}^pH^iK)_E\simeq{}^pH^i(K_E),\qquad
  ({}^{p^+}H^iK)_E\simeq{}^pH^i(K_E).
\tag{\thechapter.6}\label{eq:rationalization-perverse-cohomology}
\]
The rationalized \(p\)- and \(p^+\)-truncation triangles are canonically
identified.
\end{theorem}

\begin{proof}
For \(p\), flatness commutes with stalks, costalks, and ordinary cohomology,
so the defining vanishing inequalities are preserved.  For \(p^+\), use
\eqref{eq:pplus-tilt-formulas}: the possible \({}^pH^1\) is torsion and
vanishes after tensoring with \(E\), while the torsion-free condition on
\({}^pH^0\) introduces no additional rational cohomology.  Thus the two
integral aisles and coaisles have the same rational images.  A \(t\)-exact
functor commutes with cohomology and truncation, giving
\eqref{eq:rationalization-perverse-cohomology}.
\end{proof}

\begin{proposition}[Duality and the six operations after inverting
\(\varpi\)]
\label{prop:rationalization-six-duality}
For constructible finite-Tor objects, the canonical comparison maps give
\[
  (\D_{X,\mathcal O}K)_E\simeq\D_{X,E}(K_E).
\tag{\thechapter.7}\label{eq:rationalization-duality}
\]
Likewise rationalization commutes with \(f^*,Rf_!,Rf_*,f^!\), tensor
product, and internal Hom whenever those functors are defined on the
constructible categories in question.
\end{proposition}

\begin{proof}
This is the coefficient-change theorem for constructible complete
coefficients, not a formal claim about commuting localization with an
inverse limit.  The equivalence between constructible \(E\)-complexes and
the localization of constructible \(\mathcal O\)-complexes, together with
its tensor and six-operation compatibilities, is
\cite[Proposition~6.8.14 and Remark~6.8.15]{BhattScholzeProEtale}; see also
\cite[Theorem~7.7]{HRSConstructible}.  These comparison theorems give the
statements for \(f^*,Rf_!,Rf_*\), tensor product, and internal Hom.
The \(f^!\)-comparison follows from the adjunction with \(Rf_!\), equivalently
from the duality formula on dualizable constructible objects.  Applying the
internal-Hom comparison to the dualizing object gives
\eqref{eq:rationalization-duality}.
\end{proof}

\begin{warning}
The converse to \cref{thm:rationalization-perverse} is false without an
integral boundedness condition: rationalization cannot see a torsion object
or decide whether it lies in perverse degree \(0\) or \(1\).  Consequently
one must never prove an integral \(t\)-exactness statement merely by
tensoring with \(E\).
\end{warning}

\section{Perverse lattices}

\begin{definition}[Perverse lattice]
\label{def:perverse-lattice}
Let \(P_E\in\Perv(X,E)\).  A \(p\)-perverse lattice in \(P_E\) is an object
\(P_{\mathcal O}\in\Perv_p(X,\mathcal O)\) such that multiplication by
\(\varpi\) is monic and \(P_{\mathcal O,E}\simeq P_E\).  Two lattices are
commensurable if they become comparable after multiplying by powers of
\(\varpi\).
\end{definition}

\begin{lemma}[A torsion-free \(p\)-object lies in both hearts]
\label{lem:lattice-in-both-hearts}
If \(P_{\mathcal O}\) is a perverse lattice, then it belongs to both
\(\Perv_p(X,\mathcal O)\) and \(\Perv_{p^+}(X,\mathcal O)\).
\end{lemma}

\begin{proof}
Its \(p\)-cohomology is concentrated in degree \(0\), and that degree is
torsion-free.  Both conditions in \eqref{eq:pplus-tilt-formulas} follow.
\end{proof}

If \(U\) is smooth of pure dimension \(d\) and \(L_E\) is a lisse
\(E\)-sheaf, the image of the geometric fundamental group is compact.  After
a finite extension of \(E\), it stabilizes a finite free
\(\mathcal O\)-lattice \(L_{\mathcal O}\).  Then
\(L_{\mathcal O}[d]\) is a perverse lattice in \(L_E[d]\), common to the two
hearts.  The choice is not canonical, but any two choices are commensurable.

\section{Two integral intermediate extensions, one rational extension}

Let \(j:U\hookrightarrow X\) be an open immersion and let
\(P_{\mathcal O}\) lie in both integral hearts.  Form the images in the two
different abelian categories:
\begin{equation}
\begin{aligned}
 {}^p j_{!*}P_{\mathcal O}
   &:=\im_{\Perv_p}\bigl({}^pH^0Rj_!P_{\mathcal O}\to
                         {}^pH^0Rj_*P_{\mathcal O}\bigr),\\
 {}^{p^+}j_{!*}P_{\mathcal O}
   &:=\im_{\Perv_{p^+}}\bigl({}^{p^+}H^0Rj_!P_{\mathcal O}\to
                             {}^{p^+}H^0Rj_*P_{\mathcal O}\bigr).
\end{aligned}
\label{eq:two-integral-jmiddle}
\end{equation}

\begin{theorem}[Integral and rational middle extensions]
\label{thm:two-integral-intermediate-extensions}
The two objects in \eqref{eq:two-integral-jmiddle} need not be isomorphic
over \(\mathcal O\).  Their images in the Verdier localization by
\(\varpi\)-primary torsion are canonically isomorphic; equivalently, their
rationalizations are canonically isomorphic:
\[
  ({}^p j_{!*}P_{\mathcal O})_E
  \simeq j_{!*}(P_E)\simeq
  ({}^{p^+}j_{!*}P_{\mathcal O})_E.
\tag{\thechapter.9}\label{eq:two-jmiddle-rationalize}
\]
There is in general no canonical integral comparison morphism between the
two lattices.  If one chooses an integral comparison whose rationalization
is the displayed isomorphism, the cohomology objects of its cone are
\(\varpi\)-primary torsion.  Verdier duality exchanges the two constructions:
\[
  \D_{X,\mathcal O}({}^p j_{!*}P_{\mathcal O})\simeq
  {}^{p^+}j_{!*}(\D_{U,\mathcal O}P_{\mathcal O}).
\tag{\thechapter.10}\label{eq:integral-jmiddle-duality}
\]
\end{theorem}

\begin{proof}
By \cref{thm:rationalization-perverse}, rationalization is exact on either
heart and sends the two maps defining \eqref{eq:two-integral-jmiddle} to the
same morphism
\({}^pH^0Rj_!P_E\to{}^pH^0Rj_*P_E\).  Exact functors preserve images, which
proves \eqref{eq:two-jmiddle-rationalize} and the canonical identification in
the Verdier quotient.  A chosen integral comparison becomes an isomorphism
after tensoring with \(E\); the localization criterion therefore says
precisely that every cohomology object of its cone is torsion.

Duality exchanges \(Rj_!\) and \(Rj_*\), reverses the defining arrow, and by
\cref{thm:integral-verdier-exchange} exchanges the two hearts.  The image of
an arrow in one heart is sent to the image of the reversed dual arrow in the
other heart, giving \eqref{eq:integral-jmiddle-duality}.  Equivalently, use
the minimal-extension characterization: no nonzero boundary subobject or
quotient is carried by duality to the corresponding condition in the other
heart.  Compare \cite[\S1.4.22--1.4.25 and \S3.3.4]{BBD}.
\end{proof}

\begin{warning}[There is no canonical integral intersection complex]
For \(X\) irreducible with smooth dense \(j:U\hookrightarrow X\), a free
lattice \(L_{\mathcal O}\subset L_E\) produces
\[
  {}^p\IC_X(L_{\mathcal O})={}^p j_{!*}L_{\mathcal O}[d],\qquad
  {}^{p^+}\IC_X(L_{\mathcal O})={}
  ^{p^+}j_{!*}L_{\mathcal O}[d].
\]
Both rationalize to \(\IC_X(L_E)\), but they may differ by torsion.  Integral
duality exchanges them.  The rational intersection complex is canonical;
an integral lattice in it is additional data.
\end{warning}

\section{Reduction modulo \texorpdfstring{\(\varpi\)}{the uniformizer}}

Unlike rationalization, reduction is not flat.  The triangle
\[
  K\xrightarrow{\ \varpi\ }K\longrightarrow K_k\longrightarrow K[1]
\tag{\thechapter.11}\label{eq:derived-reduction-triangle}
\]
gives, for \(P\in\Perv_p(X,\mathcal O)\), an exact sequence
\[
  0\longrightarrow{}^pH^{-1}(P_k)\longrightarrow P
  \xrightarrow{\ \varpi\ }P\longrightarrow{}^pH^0(P_k)\longrightarrow0.
\tag{\thechapter.12}\label{eq:perverse-reduction-tor}
\]
Thus
\[
 {}^pH^{-1}(P_k)=P[\varpi],\qquad
 {}^pH^0(P_k)=P/\varpi P.
\]
Reduction of a perverse lattice is perverse because the first term vanishes.
Reduction of an arbitrary integral perverse sheaf may occupy two adjacent
perverse degrees.  This is the modular shadow of the \(p/p^+\) distinction.

\begin{proposition}[Safe lattice route to rational statements]
\label{prop:safe-lattice-route}
For the Weil proof, the following order is valid:
\begin{enumerate}
\item construct the six operations and Verdier duality over \(\mathcal O\);
\item select lattices only where compactness or a geometric construction
      provides them;
\item rationalize to \(E\), where the middle perverse structure is
      self-dual and weights are defined;
\item prove purity and the Weil bounds over \(E\);
\item return to the chosen lattices for integrality arguments.
\end{enumerate}
No step requires choosing one of the two integral intermediate extensions as
the canonical rational one.
\end{proposition}

\begin{proof}
Steps 1--3 are justified by
\cref{thm:integral-verdier-exchange,thm:rationalization-perverse,prop:rationalization-six-duality}.
Weights concern eigenvalues in \(\Qbarl\), so they are rational notions.
Both integral intermediate extensions have the same rationalization by
\cref{thm:two-integral-intermediate-extensions}.  Finally, a
Frobenius-stable lattice makes the characteristic polynomial integral by
Cayley--Hamilton; that later algebraic argument depends on the lattice, not
on a false integral self-duality assertion.
\end{proof}

\begin{auditbox}[title={Dependency firewall for integral coefficients}]
The following implications are forbidden:
\[
\begin{aligned}
 \text{pro-etale}&\Longrightarrow\text{integral self-duality of \(p\)},\\
 K_E\text{ perverse}&\Longrightarrow K\text{ perverse},\\
 P\text{ integral perverse}&\Longrightarrow P_k\text{ perverse},\\
 {}^p j_{!*}P&={} ^{p^+}j_{!*}P.
\end{aligned}
\]
The valid replacements are respectively
\cref{thm:integral-verdier-exchange,thm:rationalization-perverse},
\eqref{eq:perverse-reduction-tor}, and
\eqref{eq:two-jmiddle-rationalize}.
\end{auditbox}

\section{Formalization interface}

\begin{formalization}
A reusable implementation should expose the following objects separately.
\begin{enumerate}
\item A coefficient record \((\mathcal O,\varpi,E,k_{\mathcal O})\), with
      flat localization \(\mathcal O\to E\) and a length-one resolution of
      \(k_{\mathcal O}\).
\item The torsion pair \((\mathsf T_X,\mathsf F_X)\) in the \(p\)-heart.
\item A generic right-tilt construction producing \(p^+\) and proving
      \eqref{eq:pplus-tilt-formulas}.
\item An anti-\(t\)-exactness theorem for Verdier duality between two named
      \(t\)-structures, not an overloaded claim that it preserves one.
\item Two \(t\)-exactness instances for rationalization, both landing in the
      same rational perverse structure.
\item Separate definitions of \({}^p j_{!*}\) and
      \({}^{p^+}j_{!*}\), followed by a theorem identifying only their
      rationalizations.
\end{enumerate}
Regression tests should include the point object
\(\mathcal O/\varpi\): its dual must occur in degree \(1\), and its derived
reduction must exhibit the Tor term in degree \(-1\).
\end{formalization}

